\begin{enumerate}[label=\thesubsection.\arabic*, ref=\thesubsection.\theenumi,itemsep=0pt]

\item Dot Product of Two Vectors
\begin{itemize}
\item Matrix and Summation Forms:
For two $n$-dimensional vectors $\vec{a}$ and $\vec{b}$,
\begin{align}
\vec{a}^{\top}\vec{b} &= \sum_{i=1}^{n}a_i b_i\label{eq:dot_product}
\end{align}
\end{itemize}

\item Argument of Product of Complex Numbers
\begin{itemize}
\item Euler Form:
For $z_1=r_1e^{i\theta_1}$ and $z_2=r_2e^{i\theta_2}$,
\begin{align}
z_1 &= r_1e^{i\theta_1}, \qquad z_2=r_2e^{i\theta_2}\label{eq:complex_euler_forms}\\
z_1z_2 &= r_1r_2e^{i(\theta_1+\theta_2)}\label{eq:complex_product}\\
\arg(z_1z_2) &= \theta_1+\theta_2 \pmod{2\pi}\label{eq:arg_product}\\
\theta_1+\theta_2 &= \arg(z_1z_2)+2n\pi,\qquad n\in\mathbb{Z}\label{eq:arg_product_general}
\end{align}
\end{itemize}

\item Mean Residence Time from RTD
\begin{itemize}
\item Mean Residence Time:
For a residence time distribution $E(t)$,
\begin{align}
\bar{t} &= \frac{\int_0^\infty tE(t)\,dt}{\int_0^\infty E(t)\,dt}\label{eq:mean_rtd_ratio}\\
\int_0^\infty E(t)\,dt &= 1\label{eq:rtd_normalization}\\
\bar{t} &= \int_0^\infty tE(t)\,dt\label{eq:mean_rtd_integral}
\end{align}
\end{itemize}

\item Permutations and Combinations
\begin{itemize}
\item Number of Ways of Choosing $r$ Items from $n$ Items:
\begin{align}
{}^nC_r &= \binom{n}{r}=\frac{n!}{r!(n-r)!}\label{eq:combination_formula}
\end{align}
\end{itemize}

\item Expectation of a Function
\begin{itemize}
\item Continuous Random Variable:
For a continuous random variable $X$ with probability density function $f_X(x)$,
\begin{align}
\mathbb{E}[g(X)] &= \int_{-\infty}^{\infty}g(x)f_X(x)\,dx\label{eq:expectation_continuous}\\
\mathbb{E}[X] &= \int_{-\infty}^{\infty}xf_X(x)\,dx\label{eq:expectation_x}
\end{align}
\item Discrete Random Variable:
\begin{align}
\mathbb{E}[g(X)] &= \sum_x g(x)\Pr(X=x)\label{eq:expectation_discrete}
\end{align}
\end{itemize}

\item Newton-Raphson Method
\begin{itemize}
\item First-Order Recurrence:
For a differentiable function $f(x)$, the root is iteratively approximated using the Newton-Raphson recurrence
\begin{align}
x_{n+1} &= x_n-\frac{f(x_n)}{f'(x_n)}\label{eq:newton_raphson}
\end{align}
\end{itemize}

\item Eigenvalues, Eigenvectors and Trace
\begin{itemize}
\item Definition:
For a square matrix $\vec{A}$,
\begin{align}
\vec{A}\vec{x} &= \lambda\vec{x}, \qquad \vec{x}\ne\vec{0}\label{eq:eigen_definition}
\end{align}
where $\lambda$ is an eigenvalue and $\vec{x}$ is the corresponding eigenvector.

\item Characteristic Equation:
\begin{align}
\det(\vec{A}-\lambda\vec{I}) &= 0\label{eq:characteristic_equation}
\end{align}

\item Relation with Trace:
If $\lambda_1,\lambda_2,\ldots,\lambda_n$ are the eigenvalues of $\vec{A}\in\mathbb{R}^{n\times n}$,
\begin{align}
\operatorname{tr}(\vec{A}) &= \sum_{i=1}^{n}\lambda_i
= \lambda_1+\lambda_2+\cdots+\lambda_n\label{eq:eigen_trace}
\end{align}
Also, from the eigenvalue product relation,
\begin{align}
\det(\vec{A}) &= \prod_{i=1}^{n}\lambda_i\label{eq:eigen_determinant}
\end{align}
\end{itemize}

\item Trapezoidal Method for Plotting Differential Equations
\begin{itemize}
\item Differential Equation:
\begin{align}
\frac{dy}{dx} &= f(x,y)
\end{align}

\item Trapezoidal Update:
\begin{align}
y_{n+1}
&= y_n+\frac{h}{2}
\left[f(x_n,y_n)+f(x_{n+1},y_{n+1})\right]
\label{eq:trapezoidal_method}
\end{align}

\item The method approximates the solution by using the average of the slopes at two consecutive points. The calculated points can then be used to plot the numerical solution curve.
\end{itemize}

\item Jacobian Matrix and Divergence
\begin{itemize}
\item Jacobian Matrix:
For the vector field
$\vec{F}=\myvec{F_1 & F_2 & \cdots & F_n}^{\top}$ with variables
$\vec{x}=\myvec{x_1 & x_2 & \cdots & x_n}^{\top}$,
\begin{align}
J_{\vec{F}} &=
\myvec{
\frac{\partial F_1}{\partial x_1} & \frac{\partial F_1}{\partial x_2} & \cdots & \frac{\partial F_1}{\partial x_n}\\
\frac{\partial F_2}{\partial x_1} & \frac{\partial F_2}{\partial x_2} & \cdots & \frac{\partial F_2}{\partial x_n}\\
\vdots & \vdots & \ddots & \vdots\\
\frac{\partial F_n}{\partial x_1} & \frac{\partial F_n}{\partial x_2} & \cdots & \frac{\partial F_n}{\partial x_n}
}\label{eq:jacobian}
\end{align}

\item Divergence:
\begin{align}
\nabla\cdot\vec{F} &= \sum_{i=1}^{n}\frac{\partial F_i}{\partial x_i}\label{eq:divergence}
\end{align}

\item Relation Between Jacobian and Divergence:
\begin{align}
\nabla\cdot\vec{F} &= \operatorname{tr}(J_{\vec{F}})\label{eq:divergence_trace}
\end{align}
\end{itemize}

\item Dead Time and Time Delay in the Laplace Domain
\begin{itemize}
\item Time Delay:
A pure time delay of $\theta$ in the time domain is represented in the Laplace domain by
\begin{align}
e^{-\theta s}
\end{align}

\item Delayed Signal:
If a signal $x(t)$ is delayed by $\theta$, then
\begin{align}
x_d(t) &= x(t-\theta)u(t-\theta)\\
X_d(s) &= e^{-\theta s}X(s)
\end{align}

\item Significance in a Transfer Function:
For a process transfer function
\begin{align}
G(s) &= G_0(s)e^{-\theta s},
\end{align}
the factor $e^{-\theta s}$ represents a pure time delay or dead time of $\theta$ time units, while $G_0(s)$ represents the process dynamics without the delay.

\item The dead time shifts the process response along the time axis by $\theta$ without changing the steady-state gain. For a first-order plus dead-time model,
\begin{align}
G_{\mathrm{FOPDT}}(s) &= \frac{Ke^{-\theta s}}{\tau s+1}
\end{align}
where $K$ is the process gain, $\tau$ is the time constant, and $\theta$ is the dead time. This is the same FOPDT form given in \eqref{eq:fopdt_model}.
\end{itemize}

\item Maximum Slope Method for FOPDT Approximation
\begin{itemize}
\item First-Order Plus Dead-Time Model:
\begin{align}
G(s) &= \frac{Ke^{-\theta s}}{\tau s+1}\label{eq:fopdt_model}
\end{align}

\item Process Gain (from the FOPDT model in \eqref{eq:fopdt_model}):
\begin{align}
K &= \frac{\Delta y}{\Delta u}\label{eq:process_gain}
\end{align}

\item Maximum Slope:
\begin{align}
R &= \max\left(\frac{dy}{dt}\right)\label{eq:maximum_slope}
\end{align}

\item Dead Time:
Using the maximum-slope construction in \eqref{eq:maximum_slope}, if $(t_{\mathrm{inflection}},y_{\mathrm{inflection}})$ is the inflection point,
\begin{align}
t_i &= t_{\mathrm{inflection}}
-\frac{y_{\mathrm{inflection}}-y_0}{R}\label{eq:dead_time_tangent}\\
\theta &= t_i-t_0\label{eq:dead_time}
\end{align}

\item Time Constant:
Using the process gain and maximum slope definitions in \eqref{eq:process_gain} and \eqref{eq:maximum_slope},
\begin{align}
\tau &= \frac{K\Delta u}{R}\label{eq:fopdt_time_constant}
\end{align}
\end{itemize}

\item Transfer Function and Feedback Loop
\begin{itemize}
\item Transfer Function:
\begin{align}
G(s) &= \frac{Y(s)}{X(s)}\label{eq:transfer_function}\\
Y(s) &= G(s)X(s)\label{eq:transfer_relation}
\end{align}

\item Negative Feedback:
For forward transfer function $G(s)$ and feedback transfer function $H(s)$,
\begin{align}
E(s) &= X(s)-H(s)Y(s)\label{eq:feedback_error}\\
Y(s) &= G(s)E(s)\label{eq:feedback_output}
\end{align}

\item Closed-Loop Transfer Function:
Using the relations in \eqref{eq:feedback_error} and \eqref{eq:feedback_output},
\begin{align}
Y(s) &= G(s)\left[X(s)-H(s)Y(s)\right]\label{eq:closed_loop_output}\\
Y(s)\left[1+G(s)H(s)\right] &= G(s)X(s)\label{eq:closed_loop_relation}\\
\frac{Y(s)}{X(s)} &= \frac{G(s)}{1+G(s)H(s)}\label{eq:closed_loop_tf}
\end{align}

\item Unity Feedback:
Setting the feedback transfer function to unity in \eqref{eq:closed_loop_tf},
\begin{align}
H(s) &= 1\label{eq:unity_feedback}\\
\frac{Y(s)}{X(s)} &= \frac{G(s)}{1+G(s)}\label{eq:unity_closed_loop}
\end{align}
\end{itemize}

\item Routh-Hurwitz Stability Criterion
\begin{itemize}
\item Characteristic Equation:
For a characteristic polynomial
\begin{align}
a_ns^n+a_{n-1}s^{n-1}+\cdots+a_1s+a_0 &= 0\label{eq:routh_characteristic}
\end{align}

\item Routh Array:
\begin{align}
\myvec{
a_n & a_{n-2} & a_{n-4} & \cdots\\
a_{n-1} & a_{n-3} & a_{n-5} & \cdots\\
b_1 & b_2 & b_3 & \cdots\\
c_1 & c_2 & c_3 & \cdots
} &\label{eq:routh_array}
\end{align}

\item First Elements:
From the Routh array structure in \eqref{eq:routh_array}, the first elements are
\begin{align}
b_1 &= \frac{a_{n-1}a_{n-2}-a_na_{n-3}}{a_{n-1}}\label{eq:routh_b1}\\
b_2 &= \frac{a_{n-1}a_{n-4}-a_na_{n-5}}{a_{n-1}}\label{eq:routh_b2}
\end{align}

\item Number of Unstable Poles:
\begin{align}
N_{\mathrm{RHP}} &= \text{number of sign changes in the first column}\label{eq:routh_unstable_poles}
\end{align}

\item Stability Conditions:
\begin{align}
\text{Stable system} &\iff \text{no sign changes in the first column}\label{eq:routh_stable}\\
\text{Unstable system} &\iff \text{one or more sign changes in the first column}\label{eq:routh_unstable}\\
\text{Number of unstable poles} &= N_{\mathrm{RHP}}\label{eq:routh_rhp_count}
\end{align}
\end{itemize}

\end{enumerate}
